\subsection{FIRST PROPERTIES OF COXETER GROUPS}

Recall that from now on we assume that the elements of S are of order 2.

DEFINITION 3. (W,S) is said to be a Coxeter system if it satisfies the following condition:

\begin{enumerate}
\def\labelenumi{(\Alph{enumi})}
\setcounter{enumi}{2}
\tightlist
\item
  For \(s, s'\) in \(S\) , let \(m(s, s')\) be the order of \(ss'\) and let \(I\) be the set of pairs \((s, s')\) such that \(m(s, s')\) is finite. The generating set \(S\) and the relations \((ss')^{m(s, s')} = 1\) for \((s, s')\) in \(I\) form a presentation \(^{1}\) of the group \(W\) .
\end{enumerate}

When (W,S) is a Coxeter system, one also says, by abuse of language, that W is a Coxeter group.

Examples. 1) Let \(m\) be an integer \(\geqslant 2\) or \(\infty\) and let \(W\) be a group defined by a set of generators \(S = \{s, s'\}\) and the relations \(s^2 = s'^2 = 1\) when \(m = \infty\) , \(s^2 = s'^2 = (ss')^m = 1\) when \(m\) is finite. Consider on the other hand the dihedral group \(D_m\) (no. 2, Example) and the elements \(\rho\) and \(\rho'\) of \(D_m\) defined by (7). Since \(\rho^2 = \rho'^2 = 1\) and \((\rho \rho')^m = 1\) when \(m\) is finite, there exists a unique homomorphism \(f\) from \(W\) to \(D_m\) such that \(f(s) = \rho\) and \(f(s') = \rho'\) . Since \(\rho \rho'\) is of order \(m\) , it follows that \(ss'\) is itself of order \(m\) . Hence, \((W, S)\) is a Coxeter system, \(W\) is a dihedral group of order \(2m\) and \(f\) is an isomorphism (Prop. 2).

By transport of structure, it follows that every dihedral group is a Coxeter group.

\begin{enumerate}
\def\labelenumi{\arabic{enumi})}
\setcounter{enumi}{1}
\item
  Let \(\mathfrak{S}_n\) be the symmetric group of degree \(n\) , with \(n \geqslant 2\) . Let \(s_i\) be the transposition of \(i\) and \(i + 1\) for \(1 \leqslant i < n\) , and let \(S = \{s_1, \ldots, s_{n-1}\}\) . One can show (§2, no.4, Example and §1, Exerc. 4) that \((\mathfrak{S}_n, S)\) is a Coxeter system.
\item
  For the classification of finite Coxeter groups, cf.~Chap. 4, § 4.
\end{enumerate}

Remark. Suppose that \((W,S)\) is a Coxeter system. There exists a homomorphism \(\varepsilon\) from W to the group \(\{1,-1\}\) characterized by \(\varepsilon(s)=-1\) for all \(s\in S\) . We call \(\varepsilon(w)\) the signature of w; it is equal to \((-1)^{l(w)}\) . The formula \(\varepsilon(ww')=\varepsilon(w)\varepsilon(w')\) thus translates into \(l(ww')\equiv l(w)+l(w')\) mod. 2.

PROPOSITION 3. Assume that (W, S) is a Coxeter system. Then, two elements s and s' of S are conjugate \(^{2}\) in W if and only if the following condition is satisfied:

\begin{enumerate}
\def\labelenumi{(\Roman{enumi})}
\tightlist
\item
  There exists a finite sequence \((s_1, \ldots, s_q)\) of elements of \(S\) such that \(s_1 = s, s_q = s'\) and \(s_j s_{j+1}\) is of finite odd order for \(1 \leqslant j < q\) .
\end{enumerate}

Let \(s\) and \(s'\) in \(S\) be such that \(p = ss'\) is of finite order \(2n + 1\) . By (9), \(sp^{-n} = p^n s\) hence

\[
p ^ {n} s p ^ {- n} = p ^ {n} p ^ {n} s = p ^ {- 1} s = s ^ {\prime} s s = s ^ {\prime},\tag{10}
\]

and \(s'\) is conjugate to \(s\) .

For all s in S, let \(A_{s}\) be the set of \(s' \in S\) satisfying (I). With the hypotheses in (I), the elements \(s_{j}\) and \(s_{j+1}\) are conjugate for \(1 \leqslant j < q\) by the above, hence every element \(s'\) of \(A_{s}\) is conjugate to s.

Let \(f\) be the map from \(S\) to \(M = \{1, -1\}\) equal to 1 on \(A_s\) and to -1 on \(S - A_s\) . Let \(s'\) and \(s''\) in \(S\) be such that \(s's''\) is of finite order \(m\) . Then \(f(s')f(s'') = 1\) if \(s'\) and \(s''\) are both in \(A_s\) or both in \(S - A_s\) . In the other case, \(f(s')f(s'') = -1\) , but \(m\) is even. Thus \((f(s')f(s''))^m = 1\) in all cases. Since \((W, S)\) is a Coxeter system, there exists a homomorphism \(g\) from \(W\) to \(M\) inducing \(f\) on \(S\) . Let \(s'\) be a conjugate of \(s\) . Since \(s\) belongs to the kernel of \(g\) , so does \(s'\) , hence \(f(s') = g(s') = 1\) and finally \(s' \in A_s\) . Q.E.D.

